% Hand-authored prose. Every numeric value is a macro from generated/macros.tex,
% which is written from results/ by experiments/make_report_tables.py. Editing a
% number here is impossible by design -- that is the point.

Table~\ref{tab:metrics} reports out-of-sample performance for all \nTrials{}
configurations over \oosPeriods{} weeks. The equally weighted benchmark returned
\benchReturn{} annually at a Sharpe ratio of \benchSharpe{}, with a maximum
drawdown of \benchDrawdown{} and average turnover of \benchTurnover{}.

The best configuration was \bestName{}, at a Sharpe of \bestSharpe{}. That is
\bestSharpeGap{} above the benchmark. Across every configuration tried, the total
spread in Sharpe was \sharpeSpread{}.

Table~\ref{tab:significance} places those differences against their sampling
variability, and the picture changes completely. The benchmark's own 95\%
stationary-bootstrap interval runs from \benchCILow{} to \benchCIHigh{}. That is a
width of \benchCIWidth{} Sharpe units --- \ciSpreadRatio{} times the \emph{entire}
best-to-worst spread of the study (\sharpeSpread{}), and the widest interval any
strategy here attains (\widestCI{}); the median is \medianCI{}. The uncertainty
attaching to a single strategy's Sharpe ratio is therefore several times larger
than every difference the experiment set out to measure. Of the \nStrategies{}
optimised configurations,
\nBeatingBenchmark{} separated from the benchmark at the 5\% level under the
Ledoit--Wolf studentised bootstrap; the best configuration's difference carried a
$p$-value of \bestPValue{}.

Figure~\ref{fig:intervals} makes the point visually. The ordering of the point
estimates is stable enough to tabulate and far too noisy to interpret: the
intervals overlap almost completely, and a reader shown only the ranking would
draw conclusions the data do not support.

\paragraph{Reading the factorial.} Figure~\ref{fig:grid} separates the two
dimensions. Holding the allocator fixed at convex optimisation and varying the
belief across sample moments, Ledoit--Wolf shrinkage, Bayes--Stein, the
Normal--Inverse--Wishart predictive and Black--Litterman moves the Sharpe ratio
very little. Holding the belief fixed and switching the search method from convex
optimisation to Gaussian-process optimisation --- the like-for-like comparison, on
an identical objective over an identical feasible set --- changes mean Sharpe by
\optimiserGap{}, from \convexMeanSharpe{} to \gpMeanSharpe{}.

The two allocators do differ sharply in one respect that the Sharpe ratio hides.
Mean turnover under convex optimisation was \convexTurnover{}; under
Gaussian-process search it was \gpTurnover{}. The stochastic surrogate does not
return to the same weights when the objective barely changes between rebalances,
so it churns the portfolio without a corresponding gain.

The minimum-variance allocators are a separate and informative case. They ignore
$\mu$ entirely, and they achieve materially lower volatility and materially lower
returns than the maximum-Sharpe allocators. Their weaker Sharpe ratios should not
be read as a failure of the estimators feeding them: they are solving a different
problem, and their inclusion establishes how much of the maximum-Sharpe result
depends on the expected-return estimates at all.

\paragraph{Deflation.} The Deflated Sharpe Ratio column of
Table~\ref{tab:significance} adjusts each result for the fact that \nTrials{}
configurations were tried and the best one reported. Because every configuration
run in this study is counted --- including those that lost --- the deflation is
honest rather than nominal. It is worth stating plainly that had a single
configuration been run, reported, and the others discarded, the same numbers would
have looked considerably stronger without being any more true.
